\section{Implementações Conhecidas}

Trabalho substancial anterior existe sobre implementações do Modelo de Atores para criar sistemas distribuídos. Bibliotecas e frameworks principais como Akka (Java) e Orleans (C\#) fornecem suporte ao Modelo de Atores através de múltiplas linguagens. No entanto, a implementação mais proeminente e bem-sucedida é a linguagem de programação Elixir.

Na década de 1980, a empresa de telecomunicações sueca Ericsson experimentou com linguagens de programação funcional para abordar desafios de telecomunicações, desenvolvendo eventualmente Erlang. Erlang focou em concorrência, distribuição e tolerância a falhas implementando o Modelo de Atores no nível da linguagem. BEAM, a máquina abstrata que serviu como runtime para programas Erlang, representou a maior conquista da Ericsson neste domínio. (Na época, o termo "máquina virtual" ainda não era aplicado a tais sistemas.) No entanto, a sintaxe do Erlang era considerada complexa, criando uma barreira para adoção mais ampla.

No início dos anos 2010, José Valim estava desenvolvendo soluções para sistemas distribuídos, os mesmos desafios que a Ericsson havia abordado três décadas antes. Ao descobrir Erlang e BEAM, ele reconheceu sua eficácia, mas os achou difíceis de usar. Isso levou à criação de Elixir, que combinou o poder do BEAM e a flexibilidade do Modelo de Atores com uma sintaxe mais acessível inspirada em Ruby.

Ao contrário de implementações baseadas em bibliotecas que adicionam suporte ao Modelo de Atores a linguagens existentes, Elixir torna o Modelo de Atores o paradigma de desenvolvimento padrão. Isso é alcançado através de suporte nativo para green threads, tolerância a falhas e comunicação inter-processo remota—diferenciadores que estabeleceram Elixir como a principal plataforma do Modelo de Atores.

Esta linguagem dinâmica continua ganhando reconhecimento, classificando-se como a segunda linguagem de programação mais admirada por três anos consecutivos na Pesquisa de Desenvolvedores do Stack Overflow. Portanto, este trabalho usará Elixir como a implementação de referência e exemplo de estado da arte do Modelo de Atores.

\subsection{A Linguagem de Programação Elixir}

A filosofia de design do Elixir difere fundamentalmente das linguagens orientadas a objetos. Ao invés de classes instanciadas em objetos, Elixir usa módulos que são gerados como atores por design. Esta abordagem torna o Modelo de Atores não apenas um padrão, mas o paradigma central da linguagem.

Os atores no Elixir são representados por processos Erlang (doravante referidos simplesmente como processos). Estes processos leves são criados usando a função \texttt{spawn}, que recebe três parâmetros: um módulo, um nome de função daquele módulo e uma lista de argumentos. Quando chamada, spawn executa a função especificada em um novo processo e retorna um Process Identifier (PID) que serve como o endereço do ator.

A comunicação entre atores acontece através de passagem de mensagens usando duas funções primárias: \texttt{send} para despachar mensagens para um processo via seu PID, e \texttt{receive} para lidar com mensagens recebidas dentro de um processo. Este mecanismo de passagem de mensagens forma a base das capacidades concorrentes e distribuídas do Elixir.

Como Elixir constrói extensivamente no ecossistema do Erlang e executa na máquina virtual BEAM, referências tanto ao Erlang quanto ao BEAM aparecerão frequentemente ao longo desta discussão sobre as características e capacidades do Elixir.

\subsubsection{O Modelo de Atores em Elixir}

A implementação do Modelo de Atores em Elixir é notavelmente simples, requerendo esforço mínimo para criar sistemas de atores concorrentes e distribuídos. O design da linguagem torna a programação baseada em atores a abordagem natural ao invés de um padrão adicional a aprender.

\textbf{Mensagens} em Elixir são representadas por qualquer estrutura de dados válida do Elixir: átomos, strings, números, tuplas, listas, ou mesmo estruturas aninhadas complexas. Esta flexibilidade permite que desenvolvedores projetem protocolos de mensagem que melhor se adequem às necessidades de sua aplicação sem serem restringidos por formatos de mensagem rígidos.

\textbf{Endereços} são Process Identifiers do Erlang (PIDs), que são identificadores únicos automaticamente gerados pela máquina virtual BEAM quando um processo é gerado. Estes PIDs servem como o endereço do ator, permitindo roteamento de mensagens tanto em ambientes locais quanto distribuídos.

\textbf{Manipulação de mensagens} ocorre através do construct \texttt{receive}, que permite que um processo aguarde e faça correspondência de padrões com mensagens recebidas. Este mecanismo permite que atores processem diferentes tipos de mensagens com lógica de manipulação apropriada, tornando o sistema tanto robusto quanto expressivo.

\subsubsection{Noções Básicas do Elixir}

Elixir é uma linguagem de programação dinamicamente tipada, o que significa que tipos são verificados apenas em tempo de execução. Versões mais novas da linguagem estão experimentando com anotações de tipo opcionais, mas ainda estão em estágios iniciais.

Existem muitos tipos de dados nativos diferentes no Elixir:

\begin{lstlisting}[language=ruby, caption=Tipos de Dados em Elixir]
i = 123 
f = 3.14 
s = "héllo" # UTF-8 
c = 'hello' # Charlists (lista de ints)
b = true
a = :atom
\end{lstlisting}

A sintaxe especial \texttt{:somename} é usada para declarar átomos, um conceito do Erlang. Eles são constantes nomeadas por si mesmas e usadas em muitos cenários. Eles são usados para representar códigos de status, nomes de funções, tags ou mesmo representar valores de dados primitivos:

\begin{lstlisting}[language=ruby, caption=Átomos em Elixir]
true == :true
false == :false
nil == :nil
\end{lstlisting}

Quando se trata de dados compostos, Elixir tem 2 tipos principais: tuplas e listas. As primeiras têm tamanho fixo e são rapidamente indexáveis. As últimas são listas simplesmente encadeadas com comprimento dinâmico. Tuplas são muito usadas como objetos de dados estruturados simples. Os dados como um todo no Elixir são imutáveis, o que significa que você não pode alterar dados, mas apenas usá-los para produzir novos dados.

\begin{lstlisting}[language=ruby, caption=Estruturas de Dados em Elixir]
response = {:ok, "Some message"}
status = elem(response, 0) # :ok
response = put_elem(response, 1, "Another message") # {:ok, "Another message"}

data = [1, 2, 3, 4]
data = ["0" | data] # ["0", 1, 2, 3, 4]
data = data ++ [5] # ["0", 1, 2, 3, 4, 5]
\end{lstlisting}

O \texttt{=} é chamado de operador de ligação, porque não simplesmente faz atribuição, mas pode ser usado para correspondência de padrões.

\begin{lstlisting}[language=ruby, caption=Correspondência de Padrões em Elixir]
x = 1
1 = x # Isso funciona
2 = x # Isso irá falhar

{status, message} = {:fail, "the system panicked"}

IO.puts "It's a #{status}. Because #{message}"
# It's a fail. Because the system panicked
\end{lstlisting}

Mais amplamente, correspondência de padrões pode ser feita com a sintaxe \texttt{->}:

\begin{lstlisting}[language=ruby, caption=Estruturas de Controle em Elixir]
response = {:ok, "Content"}

case response do
	{:ok, content} -> IO.puts "Success with content: #{content}"
	{:err, cause} -> IO.puts "Ugh, it failed due to: #{cause}"
end
\end{lstlisting}

Em Elixir, existem algumas formas diferentes de declarar funções, e elas podem ser organizadas em namespaces com o uso de módulos:

\begin{lstlisting}[language=ruby, caption=Módulos e Funções em Elixir]
defmodule Math do
  # Função pública
  def add(a, b) do
    a + b
  end

  # Forma abreviada de uma linha
  def mul(a, b), do: a * b

  # Função privada (apenas chamável dentro do módulo)
  defp square(x), do: x * x

  # Múltiplas declarações de função com casamento de padrões
  def abs(n) when is_number(n) and n < 0, do: -n
  def abs(n) when is_number(n), do: n
end

IO.puts(Math.add(2, 3))   # 5
IO.puts(Math.mul(4, 5))   # 20
\end{lstlisting}

\subsection{Exemplo em Elixir}

Agora que você foi introduzido a alguma sintaxe do Elixir, veja um exemplo simples de como implementar um sistema usando atores em Elixir:

\begin{lstlisting}[language=ruby, caption=Exemplo de Ator Echo em Elixir]
defmodule Echo do
    def init do
        # Aguarda por uma mensagem
        receive do
            # Obtém o remetente e conteúdo da mensagem
            {sender, msg} -> 
                # Ecoa a resposta com sucesso (`:ok`)
                send(sender, {:ok, "Echoing: #{msg}"})
                # Loop
                init()
        end
    end
end

defmodule Example do
    def main do
        # Inicia o ator
        addr = spawn(Echo, :init, []) 
        # Envia uma mensagem
        send(addr, {self(), "Hello world"}) 
        
        # Aguarda a resposta
        receive do 
            {:ok, response} -> 
                # Exibe a resposta
                IO.puts(response) 
        end
    end
end
\end{lstlisting}

No exemplo acima, é definido um ator echo simples. Ele recebe uma mensagem que deve conter o endereço do remetente e alguma mensagem. Ele então envia de volta ao remetente uma mensagem com conteúdo \texttt{:ok} e a mensagem recebida prefixada com \texttt{Echoing: }. Finalmente, há uma chamada recursiva para \texttt{init} para que o ator possa lidar com mais mensagens, já que Elixir não tem loops.

No módulo \texttt{Example}, a chamada para \texttt{spawn} iniciará o ator através da função \texttt{init} do módulo \texttt{Echo} sem parâmetros. Esta função retornará o PID do processo criado que serve como o endereço do ator. A função \texttt{send} envia ao ator recém-criado uma tupla onde o primeiro elemento é o PID do processo atual e o segundo é a string \texttt{"Hello world"}.

Neste momento, a chamada para \texttt{receive} na função \texttt{init} entra em ação. Ela receberá a mensagem e retornará ao remetente uma nova tupla onde o primeiro elemento é \texttt{:ok} e o segundo é a string \texttt{"Echoing: Hello world"}. Por causa da chamada recursiva, este ator está pronto para receber mais mensagens.

A chamada para \texttt{receive} na função \texttt{main} agora lidará com a resposta do ator. Ela aguardou pacientemente por uma resposta e pode agora verificar que a resposta realmente contém um \texttt{:ok} e exibir a string \texttt{Echoing: Hello world} no terminal.

Note que esse código não é executado sequencialmente. A chamada de \texttt{send} do lado do módulo \texttt{Example} não vai aguardar o ator de echo processar a mensagem e gerar uma resposta. Assim que a mensagem é enviada o módulo \texttt{Example} fica parado na chamada \texttt{receive} aguardando uma resposta.
